\documentclass{article}
\usepackage{amsmath, amssymb, amsthm}
\title{IMC 1997, Day 1 --- Solutions}
\begin{document}
\maketitle

\section*{Problem 1}
Let $\varepsilon_n>0$ with $\varepsilon_n\to 0$. Find $\lim_{n\to\infty}\frac 1n\sum_{k=1}^n \ln(k/n + \varepsilon_n)$.

\subsection*{Solution}
$\int_0^1 \ln x\,dx = -1 = \lim \frac 1n\sum \ln(k/n)$. Since $\ln(k/n)\le \ln(k/n+\varepsilon_n)\le \ln(k/n+\varepsilon)$ for $n\ge n_0$, and $\frac 1n\sum\ln(k/n+\varepsilon)\to\int_0^1\ln(x+\varepsilon)dx\to -1$ as $\varepsilon\to 0$, sandwich gives $\lim = -1$.

\section*{Problem 2}
$\sum a_n$ converges. Do the following permuted sums converge?

(a) $a_1+a_2+a_4+a_3+a_8+a_7+a_6+a_5+a_{16}+\cdots+a_9+a_{32}+\cdots$

(b) $a_1+a_2+a_3+a_4+a_5+a_7+a_6+a_8+a_9+a_{11}+a_{13}+a_{15}+a_{10}+a_{12}+a_{14}+a_{16}+\cdots$

\subsection*{Solution}
(a) Yes. Partial sums at boundary are $L_{2^{n-1}+k} = S_{2^{n-1}} + S_{2^n} - S_{2^n-k}$, and $|L_{2^{n-1}+k}-S|<3\varepsilon$ once $2^{n-1}>n_0$.

(b) No. Take $a_n = (-1)^{n+1}/\sqrt n$. The odd-numbered odd indices $\{2^{n-2}\cdot 2+1, \ldots\}$ are grouped, and one can compute $L_{3\cdot 2^{n-2}} - S_{2^{n-1}} \ge 2^{n-2}/\sqrt{2^n} \to\infty$.

\section*{Problem 3}
Real $n\times n$ matrices $A,B$ with $A^2+B^2 = AB$. If $BA-AB$ is invertible, prove $3\mid n$.

\subsection*{Solution}
Let $\omega = e^{2\pi i/3} = -\frac 12 + i\frac{\sqrt 3}{2}$. Set $S = A+\omega B$. Then
\[S\bar S = (A+\omega B)(A+\bar\omega B) = A^2 + \bar\omega AB + \omega BA + B^2 = (1+\bar\omega)AB + \omega BA = -\omega AB + \omega BA = \omega(BA - AB),\]
using $1+\bar\omega = -\omega$. Taking determinants: $\det S\cdot \overline{\det S} = |\det S|^2\in\mathbb{R}$, while $\det(\omega(BA-AB)) = \omega^n\det(BA-AB)$ with $\det(BA-AB)\ne 0$ real. So $\omega^n\in\mathbb{R}$, forcing $3\mid n$.

\section*{Problem 4}
$1<\alpha<2$.

(a) Show $\alpha$ has a unique representation $\alpha=\prod_{i\ge 1}(1+1/n_i)$ with $n_i$ positive integers, $n_i^2\le n_{i+1}$.

(b) $\alpha$ is rational iff for some $m$ and all $k\ge m$, $n_{k+1}=n_k^2$.

\subsection*{Solution}
(a) Construct inductively: $\theta_0=\alpha$, and pick $n_k$ as the least $n$ with $1+1/n < \theta_{k-1}$. Then $1+1/n_k < \theta_{k-1} \le 1+1/(n_k-1)$, giving
\[1+1/n_{k+1} < \theta_k = \theta_{k-1}/(1+1/n_k) \le (1+1/(n_k-1))/(1+1/n_k) = 1+1/(n_k^2-1),\]
so $n_{k+1}\ge n_k^2$. Since $n_1\ge 2$, $n_k\to\infty$, so $\theta_k\to 1$ and $\alpha = \prod(1+1/n_k)$. Uniqueness: if $1+1/n_k\ge\theta_{k-1}$, then $\theta_k\le 1$, contradicting convergence to 1.

(b) Using $(1+1/M)(1+1/M^2)(1+1/M^4)\cdots = 1 + 1/(M-1)$, if the product ends with $n_{k+1}=n_k^2$, $\alpha$ is rational. Conversely, if $\alpha=p/q$, write $\theta_k=p_k/q_k$ with $p_k=p_{k-1}n_k$, $q_k=q_{k-1}(n_k+1)$. The integers $p_k-q_k$ are positive; $p_k-q_k - (p_{k-1}-q_{k-1}) = (n_k-1)p_{k-1}-n_k q_{k-1} < 0$ if $\theta_{k-1}<n_k/(n_k-1)$. So if strict inequality held forever, we get an infinite strictly decreasing sequence of positive integers. Contradiction. So eventually $\theta_{k-1}=n_k/(n_k-1)$, meaning $n_{k+1}=n_k^2$ onward.

\section*{Problem 5}
$\mathbb{R}_0^n = \{x\in\mathbb{R}^n : \sum x_i = 0\}$, $\mathbb{Z}_0^n$ its integer lattice. $\|x\|_p$ as usual.

(a) If $x\in\mathbb{R}_0^n$ with $\max x_i-\min x_i\le 1$ and $p\in[1,\infty]$, $y\in\mathbb{Z}_0^n$: $\|x\|_p\le\|x+y\|_p$.

(b) For $p\in(0,1)$, give counterexample.

\subsection*{Solution}
(a) $x\ne 0$ gives $\|x\|_\infty <1$. For $x_j\le 0$: $\max x_i\le x_j+1$; for $x_j\ge 0$: $\min x_i\ge x_j-1$.

Partition indices by signs of $y_i, x_i$ (sets $I(+,+), I(+,-), I(-,+), I(-,-), I(0)$). If $I(+,+)$ or $I(-,-)$ nonempty, $\|x+y\|_\infty\ge 1>\|x\|_\infty$. Otherwise, using the $\max/\min$ inequalities:
\[|x_j+y_j|^p \ge |y_j|-1+|x_k|^p\quad (j\in I(+,-), k\in I(-,+))\]
and symmetric. Case $|I(+,-)|\ge |I(-,+)|$: summing and simplifying gives
\[\|x+y\|_p^p - \|x\|_p^p \ge \sum|y_i| - 2|I(+,-)| = 2\sum_{I(+,-)}(y_j-1) + 2\sum_{I(+,+)}y_j \ge 0.\]

(b) For $p\in(0,1)$, $t\in(1/2,1)$ rational with $mt=l(1-t)$, $n=m+l$:
\[x_i=t\text{ (}i\le m\text{)}, x_i=t-1\text{ (else)}; y_i=-1\text{ (}i\le m\text{)}, y_{m+1}=m, y_i=0\text{ (else)}.\]
Then $\|x\|_p^p-\|x+y\|_p^p = m(t^p-(1-t)^p) + (1-t)^p - (m-1+t)^p$, positive for large $m$.

\section*{Problem 6}
$F$ a family of finite subsets of $\mathbb{N}$ with pairwise nonempty intersection.

(a) Must there exist finite $Y$ with $A\cap B\cap Y\ne\emptyset$ for all $A,B\in F$?

(b) Same question if all sets in $F$ have the same size.

\subsection*{Solution}
(a) No. $A_n=\{1,3,5,\dots,2n-1,2n\}$, $B_n=\{2,4,\dots,2n, 2n+1\}$; any two intersect but no finite $Y$ works.

(b) Yes. Prove stronger: if $F,G$ are finite-set families, every $A\in F$ meets every $B\in G$, $|A|=r$, $|B|=s$, then finite $Y$ exists.

Induct on $r+s$. Base $r=s=1$ trivial. Fix $A_0\in F, B_0\in G$. For $\emptyset\ne C\subseteq A_0\cup B_0$, let $F(C)=\{A\in F : A\cap(A_0\cup B_0)=C\}$; similarly $G(D)$. For each pair $(C,D)$: if $C\cap D\ne\emptyset$, done (use $Y_{C,D}=C$); else any $A\in F(C), B\in G(D)$ meet outside $A_0\cup B_0$, so $\tilde F(C)=\{A\setminus C\}$, $\tilde G(D)=\{B\setminus D\}$ satisfy condition with smaller $r,s$. Union of all $Y_{C,D}$ works.

\end{document}
